A Memtable, also known as a write buffer, serves as a staging area for write operations prior to their being flushed to 
disk. Its primary function is to enhance write performance by batching the most recent writes in memory. This approach 
eliminates the need for immediate disk input/output (I/O) operations for each write request. Instead, the writes are
grouped together and flushed to disk in a single operation. On the other hand, It also enhance read performance by 
providing a fast lookup for the most recent writes, since the lookup path in LSM tree is first to the memtable and then 
to the disk. There are several existing memtable implementations with their own trade-offs in terms of performance,
memory consumption, and write amplification. Let's discuss a few existing memtable implementations and later we will
propose new memtable implementations. We will also discuss the performance metrics and modeling for the memtable
implementations.

\subsection{Existing Memtables}

\Paragraph{Skip List} 
The default Memtable implementation in many systems like RocksDB~\cite{Facebook2021} and LevelDB is the Skip List. 
This structure is built with multiple layers of linked lists, each layer containing progressively fewer elements to 
ease navigation. The lowest layer encompasses all elements in sorted order, akin to a standard linked list, while the 
upper layers introduce shortcuts that expedite element access. This probabilistic arrangement guarantees logarithmic 
$O(\log n)$ search, insertion, and deletion times by enabling quick traversal across its layers—not through binary 
search, but via a strategic, level-by-level method. Consequently, Skip Lists are exceptionally suited for achieving 
balanced read and write operations, as well as efficiently supporting fully ordered database scans.

However, the layered structure and the necessity to maintain pointers for each element increase the memory requirements 
of Skip Lists. Typically, a four-layered Skip List, with each upper layer containing approximately one-fourth of the 
elements in the layer immediately below it, is utilized. This configuration demands roughly 1.33 times more memory than 
the actual data stored.

\Paragraph{Vector} A Vector functions as a dynamic, automatically resizing array, providing contiguous storage for 
efficient iteration. Insertions are optimally performed at the Vector's end, making it particularly advantageous for 
write-intensive tasks. However, its dependence on linear search $O(n)$ for indexing renders it less effective for point 
queries and deletions, as these operations necessitate sequential traversal of the Vector. The memory overhead for a 
Vector is relatively minimal, requiring only additional space for resizing along with a modest amount of metadata. Due 
to its lower memory demands, a Vector can potentially store more data, leading to fewer disk flushes. Nonetheless, 
Vectors are not well-suited for handling range queries and point queries due to the linear search requirement.

\Paragraph{Hash Skip List} The Hash Skip List combines a hash table with Skip Lists within each bucket. Keys are 
partitioned by prefix for bucket allocation, facilitating targeted insertions and lookups through a combination of 
hashing and binary search. This structure excels in handling small range queries within specific key prefixes. 
Nonetheless, executing range scans across multiple prefixes is resource-intensive, as it necessitates copying and 
sorting, thereby impacting time and memory efficiency.

The memory overhead of the Hash Skip List is higher than that of a traditional Skip List due to the extra space needed 
for the hash table and the pointers to each Skip List. The Hash Skip List is particularly well-suited for workloads
involving a high volume of small range queries within specific key prefixes.

\Paragraph{Hash Link List} Similar to the Hash Skip List, the Hash Link List incorporates a hash table framework; 
however, it differs by linking each bucket to a single linked list. This arrangement is particularly tailored for 
efficiently handling range queries limited to narrow key prefixes, typically involving a modest number of entries. 
Despite utilizing hashing for quick bucket identification, the Hash Link List relies on linear search within each 
linked list for indexing, which may reduce its efficiency for point queries and deletions compared to other structures.

The Hash Link List presents a lower memory overhead compared to the Hash Skip List, primarily because it 
eschews the layered structure inherent to Skip Lists. Consequently, this streamlined setup makes it an excellent choice 
for scenarios characterized by frequent small range queries within specific key prefixes, especially when each prefix 
is associated with a limited dataset.

\subsection{Proposed Memtables} \lipsum[1]

\Paragraph{Trie} \lipsum[1]

\Paragraph{Tree} \lipsum[1]

\subsection{Performance Modeling \& Metrics} \lipsum[1]